\largerpage
\sscp{Cuscus \it{(Phalangeridae)}}{C-08-01}
Words for `cuscus' in SHWNG languages are given below,
several of which are from **kVndoR(a) (\srf{14-03}) or **mantəR (\srf{14-04}).

Laura Arnold provides the following information (summarised from
e-mail correspondence in July 2021) regarding cuscus in the Raja
Ampat Islands, which is relevant for interpreting the glosses she provides for
these languages (Ambel, Biga, As, Matbat, Ma{\Q}ya,
Legenyem, Salawati, Batta): \it{Spilocuscus papuensis} is endemic
to Waigeo (but not found on the other islands); according to
Wikipedia, other cuscus of Raja Ampat are \it{Phalanger orientalis} (throughout),
\it{Phalanger gymnotis} (reported on Misool only),
and \it{Spilocuscus maculatus} (Misool only).
Ambel has
%the follwoing terms: \it{tamʧám} ‘cuscus (generic)’,
%\it{(tamʧám) malélen} ‘multicoloured; k.o.
%cuscus’ -- based on the descriptions,
%this may be the ‘Waigeo spotted cuscus,
%\it{Spilocuscus papuensis}’,
\it{(tamʧám) hu}, glossed locally
as ‘kuskus tanah’ = ‘ground cuscus’,
which may be the ‘ground cuscus,
\it{Phalanger gymnotis}’ which lives in burrows
and rock holes during the day
and feeds in the trees at night,
but it is not entirely clear exactly what ‘kuskus tanah’ refers to.

\cite{Price2002} has extremely detailed zoological
information for Ambai, and includes many additional words for
kinds of cuscus (or possums) which are compounds for \it{fimuna}.\footnote{
		These Ambai phrases are:
		\it{fimuna embai aweŋ} `species of possum',
		\it{fimuna mantatowi} `Spotted Cuscus (in Normal phase)',
		\it{fimuna manggahai} `species of cuscus',
		\it{fimuna mera} `Spotted Cuscus (in Ginger phase)',
		\it{fimuna merasaroŋ} `species of cuscus', and
		\it{fimuna nu(m)baibua} `Spotted Cuscus (in White phase).
		In addition the phrase \it{wintumaŋ kamirei roroŋ} occurs for `ground cuscus'.
		This is an apparent synonym for \it{tumaŋ}.}
Similarly, \cite{PriceDonohue2009} have many words for
kinds of cuscus which are compounds with \it{kapa} `mammal'.\footnote{
		Ansus phrases involving \it{kapa} in \citet{PriceDonohue2009} are:
		\it{kapa dupapei} `spotted cuscus',
		\it{kapa dura} `red cuscus',
		\it{kapa duwua} `white cuscus',
		\it{kapa gondoi} `k.o. terrestrial cuscus',
		\it{kapa manu} `rat', and
		\it{kapa musi} `tree kangaroo'.}

\noindent{\wg\st{1.5pt}\begin{xltabular}{\linewidth}
{lllll>{\raggedright\arraybackslash}X}
\lsptoprule
%Language	&	ISO	&	Term	&	Gloss	&	Etymon	&	Source	\\	\midrule
%\endfirsthead
Language	&	ISO	&	Term					&	Gloss						&	Etymon				&	Source	\\	\midrule	\hhline{~}
\endhead	
Ambel			&	wgo	&	tamʧám;				&	cuscus;					&	\it{misc.}		&	\cite[777]{Arnold2018}, LA\footnotemark	\\
					&			&	tamʧám hu;		&	ground cuscus;	&								&	\hp{\,}	\\
					&			&	tamʧám 				&	kind of cuscus	&								&	\hp{\,}	\\	\hhline{~}
					&			&	malélen				&									&								&	\hp{\,}	\\
Biga			&	bhc	&	ta³f;					&	cuscus;					&	\it{misc.}		&	\hp{\,}	\\
					&			&	kaˈsu³t				&	ground cuscus		&	\it{misc.}		&	LA	\\
As				&	asz	&	do³;					&	cuscus;					&	**kVndoR(a)/	&	LA	\\	\hhline{~}
					&			&								&									&	**kindoR			&	\hp{\,}	\\
					&			&	le³ntati gu³	&	ground cuscus		&	\it{misc.}		&	\hp{\,}	\\
Matbat		&	xmt	&	doː (?);			&	cuscus;					&	**kVndoR(a)/	&	\cite[241]{Blust1982};	\\	\hhline{~}
					&			&								&									&	**kindoR;			&	\hp{\,}	\\
					&			&	iba¹²¹ŋ				&	cuscus					&	\it{misc.}		&	\cite{Schapper2011}	\\
Ma{\Q}ya	&	slz	&	amˈʧa¹²m			&	cuscus					&	\it{misc.}		&	LA	\\
Legenyem	&	lcc	&	tamˈca¹²m			&	cuscus					&	\it{misc.}		&	LA	\\
Salawati	&	xmx	&	yo³l;					&	cuscus;					&	\it{misc.}		&	LA	\\
					&			&	əsu¹²k				&	ground cuscus		&	\it{misc.}		&	\hp{\,}	\\
Batta			&	xmx	&	yo¹l;					&	cuscus;					&	\it{misc.}		&	LA	\\
					&			&	qəwi¹²t				&	ground cuscus		&	\it{misc.}		&	\hp{\,}	\\
Buli			&	bzq	&	do						&	small marsupial	&	**kVndoR(a)/	&	\cite[241]{Blust1982}	\\	\hhline{~}
					&			&								&									&	**kindoR			&	\hp{\,}	\\
Sawai			&	szw	&	sɛfsɛf				&	cuscus					&	\it{misc.}		&	\cite{WhislerWhisler1995b}	\\
Patani		&	ptn	&	sifsɛf				&	cuscus					&	\it{misc.}		&	Linne Iren Sjånes Rødvand (pers. comm. August 2023)	\\
Giman			&	gzn	&	do						&	cuscus					&	**kVndoR(a)/	&	\cite{Teljeur-Grimes1990}	\\	\hhline{~}
					&			&								&									&	**kindoR		&	\hp{\,}	\\
Taba	&	mky	&	kusok					&	cuscus					&	*kusor			&	\cite[105]{Collins1982a}	\\
\end{xltabular}}
\footnotetext{LA = Laura Arnold (pers. comm. June 2022).}

\noindent{\wg\st{1.5pt}\begin{xltabular}{\linewidth}
{lllll>{\raggedright\arraybackslash}X}
\lsptoprule
%Language	&	ISO	&	Term	&	Gloss	&	Etymon	&	Source	\\	\midrule
%\endfirsthead
Language	&	ISO	&	Term					&	Gloss						&	Etymon				&	Source	\\	\midrule	\hhline{~}
\endhead	
Biak			&	bhw	&	rambav				&	cuscus					&	\it{misc.}	&	\cite{Heuvel2019}	\\
Roon			&	rnn	&	maŋgadiaβi		&	opossum					&	\it{misc.}	&	SV \citeyear[248]{SmitsVoorhoeve1992b}	\\
Dusner		&	dsn	&	urɛm					&	opossum					&	\it{misc.}	&	SV \citeyear[248]{SmitsVoorhoeve1992b}	\\
Meoswar		&	mvx	&	karau					&	opossum					&	\it{misc.}	&	SV \citeyear[248]{SmitsVoorhoeve1992b}	\\
Yaur			&	jau	&	ˈubiɛ					&	opossum					&	\it{misc.}	&	SV \citeyear[249]{SmitsVoorhoeve1992b}	\\
Yerisiam	&	ire	&	átóòrà 				&	cuscus					&	**kVndoR(a)	&	\cite[69]{Kamholz2014}	\\
Yeretaur	&	gop	&	naˈgɔmbrar		&	opossum					&	\it{misc.}	&	SV \citeyear[249]{SmitsVoorhoeve1992b}	\\
Moor			&	mhz	&	aˈtiɛta/			&	opossum					&	\it{misc.}	&	SV \citeyear[249]{SmitsVoorhoeve1992b}	\\	\hhline{~}
					&			&	aˈtiat				&									&							&	\hp{\,}	\\
Warembori	&	wsa	&	maya-ro/			&	cuscus					&	**mantəR		&	\cite{Donohue1999b}	\\	\hhline{~}
					&			&	maye					&									&							&	\hp{\,}	\\
Yoke			&	yki	&	maraɣi				&	cuscus					&	**mantəR		&	\cite{Donohue1999b}	\\
Ansus			&	and	&	amu						&	cuscus					&	\it{misc.}	&	\cite{PriceDonohue2009}	\\
					%&			&	kapa					&	mammal					&	\it{misc.}	&	\hp{\,}	\\
					%&			&	kapa 					&	spotted cuscus	&							&	\hp{\,}	\\	\hhline{~}
					%&			&	dupapei;			&									&							&	\hp{\,}	\\
					%&			&	kapa dura;		&	red cuscus			&							&	\hp{\,}	\\
					%&			&	kapa duwua;		&	white cuscus		&							&	\hp{\,}	\\
					%&			&	kapa 					&	k.o. terrestrial&							&	\hp{\,}	\\	\hhline{~}
					%&			&	gondoi;				&	cuscus					&							&	\hp{\,}	\\
					%&			&	kapa manu;		&	rat							&							&	\hp{\,}	\\
					%&			&	kapa musi			&	tree kangaroo		&	**musi				&	\hp{\,}	\\
Ambai			&	amk	&	fimuna;				&	cuscus;					&	\it{misc.}	&	\cite{Price2002}	\\
					%&			&	fimuna 				&	species of			&							&	\hp{\,}	\\	\hhline{~}
					%&			&	embai					&	possum;					&							&	\hp{\,}	\\	\hhline{~}
					%&			&	aweŋ;					&									&							&	\hp{\,}	\\
					%&			&	fimuna 				&	Spotted Cuscus 	&							&	\hp{\,}	\\	\hhline{~}
					%&			&	mantatowi;		&	(Normal phase);	&							&	\hp{\,}	\\
					%&			&	fimuna 				&	species of 			&							&	\hp{\,}	\\	\hhline{~}
					%&			&	manggahai;		&	cuscus;					&							&	\hp{\,}	\\
					%&			&	fimuna 				&	Spotted Cuscus		&							&	\hp{\,}	\\	\hhline{~}
					%&			&	mera;					&	(Ginger phase);	&							&	\hp{\,}	\\
					%&			&	fimuna 				&	species of 			&							&	\hp{\,}	\\	\hhline{~}
					%&			&	merasaroŋ;		&	cuscus;					&							&	\hp{\,}	\\
					%&			&	fimuna 				&	Spotted Cuscus	&							&	\hp{\,}	\\	\hhline{~}
					%&			&	nu(m)baibua;	&	(White phase);	&							&	\hp{\,}	\\
					&			&	tumaŋ;				&	ground cuscus			&	\it{misc.}	&	\hp{\,}	\\	\hhline{~}
					&			&								&	\it{Strigocuscus}		&							&	\hp{\,}	\\	\hhline{~}	
					&			&								&	\it{gymnotis}				&							&	\hp{\,}	\\	
					&			&	waiwiraŋ 			&	common cuscus			&	\it{misc.}	&	\hp{\,}	\\	\hhline{~}
					&			&	tuwa;					&	\it{Phalanger }			&							&	\hp{\,}	\\	\hhline{~}
					&			&								&	\it{orientalis}			&							&	\hp{\,}	\\
					&			&	wintumaŋ			&	species of ringtail &	\it{misc.}	&	\hp{\,}	\\	\hhline{~}
					&			&								&	possum/cuscus				&							&	\hp{\,}	\\
					%&			&	wintumaŋ 		&	ground cuscus		&							&	\hp{\,}	\\	\hhline{~}
					%&			&	kamirei 		&									&							&	\hp{\,}	\\	\hhline{~}
					%&			&	roroŋ				&									&							&	\hp{\,}	\\	
Wandamen	&	wad	&	aniai 				&	opossum					&	\it{misc.}	&	SV \citeyear[248]{SmitsVoorhoeve1992b}	\\	%\hhline{~}
%					&			&	(kaˈkɔpa)			&									&							&	\hp{\,}	\\
Serui-Laut&	seu	&	wondoi				&	opossum					&	\it{misc.}	&	SV \citeyear[248]{SmitsVoorhoeve1992b}	\\
\lspbottomrule
\end{xltabular}}
